% Options for packages loaded elsewhere
\PassOptionsToPackage{unicode}{hyperref}
\PassOptionsToPackage{hyphens}{url}
\PassOptionsToPackage{dvipsnames,svgnames,x11names}{xcolor}
\documentclass[
  11pt]{article}
\usepackage{xcolor}
\usepackage[margin=1in]{geometry}
\usepackage{amsmath,amssymb}
% Section numbering kept at default; \section* used for References/Appendix/Checklist
\usepackage{iftex}
\ifPDFTeX
  \usepackage[T1]{fontenc}
  \usepackage[utf8]{inputenc}
  \usepackage{textcomp} % provide euro and other symbols
\else % if luatex or xetex
  \usepackage{unicode-math} % this also loads fontspec
  \defaultfontfeatures{Scale=MatchLowercase}
  \defaultfontfeatures[\rmfamily]{Ligatures=TeX,Scale=1}
\fi
\usepackage{lmodern}
\ifPDFTeX\else
  % xetex/luatex font selection
\fi
% Use upquote if available, for straight quotes in verbatim environments
\IfFileExists{upquote.sty}{\usepackage{upquote}}{}
\IfFileExists{microtype.sty}{% use microtype if available
  \usepackage[]{microtype}
  \UseMicrotypeSet[protrusion]{basicmath} % disable protrusion for tt fonts
}{}
\makeatletter
\@ifundefined{KOMAClassName}{% if non-KOMA class
  \IfFileExists{parskip.sty}{%
    \usepackage{parskip}
  }{% else
    \setlength{\parindent}{0pt}
    \setlength{\parskip}{6pt plus 2pt minus 1pt}}
}{% if KOMA class
  \KOMAoptions{parskip=half}}
\makeatother
\usepackage{longtable,booktabs,array}
\newcounter{none} % for unnumbered tables
\usepackage{calc} % for calculating minipage widths
% Correct order of tables after \paragraph or \subparagraph
\usepackage{etoolbox}
\makeatletter
\patchcmd\longtable{\par}{\if@noskipsec\mbox{}\fi\par}{}{}
\makeatother
% Allow footnotes in longtable head/foot
\IfFileExists{footnotehyper.sty}{\usepackage{footnotehyper}}{\usepackage{footnote}}
\makesavenoteenv{longtable}
\setlength{\emergencystretch}{3em} % prevent overfull lines
\providecommand{\tightlist}{%
  \setlength{\itemsep}{0pt}\setlength{\parskip}{0pt}}
\usepackage{bookmark}
\IfFileExists{xurl.sty}{\usepackage{xurl}}{} % add URL line breaks if available
\urlstyle{same}
\hypersetup{
  colorlinks=true,
  linkcolor={blue},
  filecolor={Maroon},
  citecolor={blue},
  urlcolor={blue},
  pdfcreator={LaTeX via pandoc}}

% -----------------------------------------------------------------
% arXiv submission package — C1 unified preprint v0.3
% Source: docs/sessions/C1-unified-preprint-draft-v0.2.md
% Converted via pandoc 3.9.0.2 on 2026-05-24
% Manual edits: title block reformatted to \title/\author/\maketitle;
% added arXiv-style abstract environment.
% -----------------------------------------------------------------

\usepackage{graphicx}
\usepackage{authblk}

\title{A pipeline for cross-domain validation of self-organized criticality:\\
five systems, one method}
\author{Wan Qihui}
\affil{Structural Isomorphism Project, Independent Researcher\\
\texttt{https://structural.bytedance.city}}
\date{2026-05-24 \\ \small v0.3 unified preprint (Phase 1--5 synthesis)}

\begin{document}

\maketitle

\noindent\textbf{Keywords.} self-organized criticality; cross-domain validation; power-law;
Gutenberg--Richter; Omori--Utsu; inverse cubic law; neural avalanches;
null control; universality class; reproducibility.

\begin{abstract}

A universality-class membership claim has empirical content only if a
single, fixed analysis pipeline --- applied with no per-domain tuning
--- recovers the predicted scaling signatures across systems drawn from
very different domains, \emph{and} correctly fails to find those
signatures in matched non-class data. We assemble such a pipeline
(Clauset--Shalizi--Newman 2009 maximum-likelihood power-law fitting with
Kolmogorov--Smirnov-driven \texttt{x\_min} selection, a likelihood-ratio
test against lognormal and exponential alternatives, and Omori--Utsu
temporal-decay stacking) into one shared Python package and apply it
unchanged to five independent systems: USGS tectonic earthquakes (Phase
1), S\&P 500 daily returns (Phase 2), DeFi liquidation cascades across
three protocols (Phase 3), task-active mouse-cortex neural avalanches
(Phase 4), and a set of four synthetic non-self-organized-criticality
(non-SOC) null sources (Phase 5). On real data the pipeline recovers
canonical exponents: a Gutenberg--Richter b-value of 1.084 ± 0.005 on
37,281 earthquakes above the completeness magnitude (Omori p = 0.941 ±
0.017); an inverse-cubic tail exponent α = 2.998 ± 0.041 on 9,060 S\&P
500 daily returns; tail exponents α ∈ {[}1.567, 1.684{]} across 43,065
on-chain DeFi liquidations spanning three architecturally distinct
lending protocols (Omori p ∈ {[}0.69, 0.76{]}); and, on 1.39 M
mouse-cortex spikes, a self-organized-criticality scaling relation
satisfied to within 2 \% at every bin scale (measured γ ≈ 1.10) although
the recording's specific exponents (τ ∈ {[}2.17, 3.00{]}) place it in a
task-active sub-class rather than the canonical mean-field regime. All
four synthetic non-SOC nulls (folded normal, exponential, Poisson
inter-arrival, Poisson Omori) are correctly rejected, with
power-law-vs-alternative likelihood ratios of −16 to −45 and an
Omori-fit R² ≈ 0.002 --- ruling out the trivial failure mode ``the
pipeline fits everything as a power law.'' We report three
qualifications honestly. (i) The lognormal alternative is not rejected
in every raw-tail test: in particular, for S\&P 500 the raw-tail
likelihood ratio favors lognormal (R = −6.12), and the SOC verdict for
that system rests on exponent-band agreement (α = 2.998 vs the canonical
3.0) rather than on rejecting the lognormal. (ii) The pipeline is
constructed to detect \emph{endogenous} threshold-cascade signatures and
--- as is a known property of SOC threshold-cascade methods --- is not
expected to flag externally driven crises (Phase 5 validates only the
null-rejection direction). (iii) The project's downstream cross-domain
\emph{predictions} (Layer 4) remain unverified --- their target data has
not been collected, so we can give only prior-based credible intervals,
not frequentist confidence intervals on observed data. Within those
limits, the joint result is an internally consistent, single-pipeline
cross-domain test of self-organized-criticality universality across
geophysics, equity finance, decentralized finance, and neuroscience.
\end{abstract}

\section{Introduction}\label{introduction}

Universality classes are the sharpest tool statistical physics offers
for cross-system comparison: two systems in the same class share a small
set of critical exponents that are independent of microscopic detail
{[}1, 2{]}. The concept was extended from equilibrium critical phenomena
to non-equilibrium dynamics through the theory of self-organized
criticality (SOC) of Bak, Tang, and Wiesenfeld {[}3{]}, in which slowly
driven threshold-cascade systems generically exhibit power-law
event-size distributions, Omori-like temporal relaxation, and associated
scaling relations without parameter tuning. Tectonic seismicity is the
canonical natural realization {[}3, 4{]}, and the Gutenberg--Richter and
Omori--Utsu laws {[}5, 6{]} are its most widely reproduced quantitative
signatures. Beggs and Plenz {[}7{]} opened the biological side of the
class with cortical avalanches showing P(s) ∝ s⁻³ᐟ² and P(T) ∝ T⁻².
Sornette {[}8{]} extended the picture to financial cascades.

The empirical literature contains many single-system measurements but
few cross-system comparisons that use one fixed fitting stack. Clauset,
Shalizi, and Newman {[}9{]} argued that standard estimators ---
binned-histogram slope fits, naive \texttt{x\_min} choices --- were
producing falsely confident power-law conclusions, and that canonical
examples deserved re-testing under
maximum-likelihood-plus-Kolmogorov--Smirnov estimation with explicit
comparison to alternatives. Subsequent practice tightened the floor: a
defensible power-law claim today requires a Clauset maximum-likelihood
fit with a reported \texttt{x\_min}, a likelihood-ratio test against at
least lognormal and exponential, and a null-control check. Most
cross-domain SOC studies do not meet this standard; the typical paper is
one system deep.

The Structural Isomorphism project is an attempt to make cross-domain
``same mathematical structure'' claims operational. Its layered pipeline
(i) builds a domain-agnostic catalog of candidate systems and
observables, (ii) groups them into candidate equivalence classes from
mechanism graphs, (iii) extracts shared invariants for each class, and
(iv) issues falsifiable numerical predictions. The Layer 1/2
community-discovery step found that a single self-organized-criticality
``threshold-cascade'' cluster emerged unsupervised from the project's
pair data --- the largest community in the graph, with earthquakes, DeFi
liquidations, bank runs, flash crashes, power-grid cascades, and neural
avalanches all assigned to it. The present paper is the empirical
validation step for that cluster's core members.

This paper synthesizes the project's first five validation phases into a
single preprint. The contributions are:

\begin{enumerate}
\def\labelenumi{\arabic{enumi}.}
\tightlist
\item
  \textbf{A single fixed pipeline across four real systems.} We re-fit
  power-law tails and, where applicable, Omori temporal decay on USGS
  earthquakes (Phase 1), S\&P 500 daily returns (Phase 2), three DeFi
  lending protocols (Phase 3 --- Aave V2, Compound V2, MakerDAO), and
  mouse-cortex neural avalanches (Phase 4), all through the same code
  path with no per-domain parameter tuning.
\item
  \textbf{Null robustness.} Phase 5 runs the identical pipeline on four
  synthetic non-SOC sources and verifies they are all correctly
  rejected, ruling out ``the pipeline fits everything'' as a trivial
  explanation for the positive findings.
\item
  \textbf{An honest accounting of what the pipeline does and does not
  establish} --- including the lognormal-not-always-rejected
  qualification, the endogenous-only scope, and the unverified status of
  the project's downstream predictions.
\end{enumerate}

\textbf{Scope decision: this is the focused five-system SOC paper.} The
Structural Isomorphism project has produced two manuscript lines at
different breadths. The present paper is the \emph{focused} synthesis of
the five-system SOC core (Phases 1--5: four real threshold-cascade
systems plus the synthetic null phase). A separate, broader manuscript
--- \texttt{unified-pipeline-v0.2-2026-05-13}, here referred to as the
\emph{thirteen-system sibling paper} --- applies the same methodological
framework to thirteen systems across five distinct universality-class
parents (adding preferential attachment, Motter--Lai network cascade,
Preisach hysteresis, and Scheffer fold-bifurcation phases). The roadmap
C1 item is explicitly titled ``\emph{N systems, one method}'' and scoped
to the five-system SOC core, so C1 ships as \textbf{this focused paper}.
The thirteen-system sibling is \emph{not} superseded by, nor merged
into, this paper: the two have different theses (this paper: one SOC
universality class verified deeply across four real domains plus a null;
the sibling: one \emph{methodological framework} shown to be portable
across five class families). They share the Phase 1--4 numbers, and this
paper is kept numerically consistent with the sibling. A reviewer may
decide at submission time whether to post both, or only the focused one;
that is an editorial choice, not an unresolved methodological gap. (For
traceability of v0.1→v0.2 changes see the META block at the top of this
file.)

The paper is organized as follows. Section 2 specifies the shared
pipeline. Section 3 reports the five phases. Section 4 places the four
real systems side by side. Section 5 discusses the cross-domain picture.
Section 6 states the limitations explicitly. Section 7 concludes.

\begin{center}\rule{0.5\linewidth}{0.5pt}\end{center}

\section{The shared pipeline}\label{sec:2}

The shared analysis stack is implemented as one Python package and
exposed to every phase as a small set of functions. The pipeline is
intentionally minimal: each step corresponds to a single published
estimator, the parameters are fixed across phases, and the only
domain-specific code lives in the per-phase data loaders. No phase
modifies the pipeline; no phase tunes a fitting parameter; no phase adds
a domain-specific prior.

\textbf{Implementation and provenance.} The authoritative implementation
is the standalone Python package \texttt{soc-pipeline} (version 0.1.0,
MIT licence), located at \texttt{packages/soc-pipeline/} in the project
repository. It is split into one module per analytical operation:
\texttt{fit.py} (Clauset maximum-likelihood power-law fit),
\texttt{bootstrap.py} (bootstrap confidence intervals),
\texttt{lr\_test.py} (likelihood-ratio tests), \texttt{omori.py}
(Omori--Utsu stacking and fitting), \texttt{null\_controls.py}
(synthetic null generators), \texttt{b\_value.py} (Aki
maximum-likelihood b-value), \texttt{universal\_collapse.py}, and
supporting utilities; it depends only on \texttt{numpy}, \texttt{scipy},
\texttt{pandas}, and \texttt{powerlaw}. A legacy path
\texttt{v4/lib/soc\_pipeline.py} is retained for backward compatibility,
but it is now a thin (≈75-line) deprecation shim that re-exports the
package and emits a \texttt{DeprecationWarning}; it is \emph{not} the
implementation. We note one provenance caveat the reviewer should be
aware of: the thirteen-system sibling manuscript describes the pipeline
as ``\texttt{v4/lib/soc\_pipeline.py}, 339 lines, frozen at commit
\texttt{7ee228c},'' a description that predates the extraction of the
code into the \texttt{soc-pipeline} package and is no longer literally
accurate for the current repository layout. The Phase 1--5 source papers
(2026-04-15/16) predate the package extraction as well and describe the
same \emph{logical} pipeline through per-phase analysis scripts
(e.g.~\texttt{v4/validation/soc-stockmarket/fetch\_and\_analyze.py},
\texttt{v4/validation/null-controls/generate\_and\_analyze.py}) tagged
per phase in the repository (\texttt{v4/phase1-earthquake-2026-04-15},
\texttt{v4/phase2-stockmarket-2026-04-15}, \ldots). The logical pipeline
is identical across all of these; before formal submission a reviewer
should record one canonical commit hash or release tag of the
\texttt{soc-pipeline} package against which the final numbers are
confirmed (see the Pre-submission checklist).

\subsubsection{2.1 Clauset--Shalizi--Newman maximum-likelihood power-law
fit}\label{clausetshalizinewman-maximum-likelihood-power-law-fit}

For each dataset we fit a continuous power-law p(s) ∝ s⁻ᵅ for s ≥
\texttt{x\_min} using the Clauset--Shalizi--Newman estimator {[}9{]}.
The lower cutoff \texttt{x\_min} is selected automatically by minimizing
the Kolmogorov--Smirnov distance between the empirical and fitted
cumulative distribution functions on the candidate tail; α is then
estimated by maximum likelihood (Hill-form estimator) on that tail. We
use the Alstott--Bullmore--Plenz \texttt{powerlaw} library {[}10{]} as
the canonical implementation, with the discrete-data option set only for
explicitly integer-valued data (Phase 4 avalanche sizes and durations).
For each fit we report α, the analytic (Hill-form) standard error σ(α),
the fitted \texttt{x\_min} in the domain's natural units, and the tail
size \texttt{n\_tail}.

\subsubsection{2.2 Uncertainty quantification and
bootstrap}\label{uncertainty-quantification-and-bootstrap}

We report two kinds of uncertainty, and --- importantly --- \emph{not
every phase uses the same one}; the source papers differ here, and we
state the per-phase situation explicitly rather than impose a false
uniformity.

\begin{itemize}
\tightlist
\item
  \textbf{Phase 1 (earthquakes)} reports a bootstrap confidence interval
  \emph{in addition to} the analytic Shi--Bolt b-value error: a 95 \%
  bootstrap CI on the Gutenberg--Richter b-value from \textbf{500
  resamples} with replacement of the events above M\_c, recomputing the
  Aki maximum-likelihood estimator on each resample, with a fixed random
  seed (42) for bit-reproducibility. The reported CI is {[}1.073,
  1.094{]}.
\item
  \textbf{Phases 2, 3, 4} (S\&P 500, DeFi protocols, neural avalanches):
  the corresponding source papers do \textbf{not} report a bootstrap.
  The α uncertainties quoted for those phases (e.g.~α = 2.998 ± 0.041,
  the DeFi α ± 0.010--0.016, the Phase 4 per-bin τ errors) are the
  \textbf{Clauset/Hill analytic standard error} returned by the
  \texttt{powerlaw} library, σ(α) ≈ (α − 1)/√n\_tail, not bootstrap
  intervals. This is the standard analytic error for a Clauset
  maximum-likelihood fit.
\end{itemize}

For completeness we note that the \texttt{soc-pipeline} package does
provide a \texttt{bootstrap\_ci} routine (default 200 resamples, ≥200
recommended), and the thirteen-system sibling manuscript additionally
reports 100-resample bootstrap CIs for its consolidated runs. But the
Phase 1--5 numbers quoted in \emph{this} paper are exactly as in the
five source papers: bootstrap for Phase 1's b-value (500 resamples),
analytic Hill error elsewhere. Appendix A, Table A2 lists the per-phase
choice. A reviewer who wishes to report a uniform bootstrap CI for all
four real phases would need to re-run Phases 2--4 through the package's
\texttt{bootstrap\_ci}; that is a possible methodological upgrade, not a
correction of an error.

\subsubsection{2.3 Likelihood-ratio tests against
alternatives}\label{likelihood-ratio-tests-against-alternatives}

For each fit we compute the Clauset--Shalizi--Newman normalized
log-likelihood ratio R against two alternatives --- lognormal and
exponential --- with associated (Vuong-style) p-values {[}9{]}. In the
Clauset convention, \textbf{a positive R favors the power-law and a
negative R favors the alternative}; p \textless{} 0.05 indicates the
preference is statistically distinguishable. Rejection of exponential is
necessary but not sufficient for a power-law claim; the harder test is
against lognormal, which can mimic a power-law tail over a finite
dynamic range. Clauset et al.~{[}9{]} caution that the likelihood-ratio
test has limited power for small tails (n\_tail \textless{}
\textasciitilde50), in which case ``inconclusive'' must not be read as
evidence for either model. The sign convention matters: §3.2 and §6.1
below turn on it.

\subsubsection{2.4 Omori--Utsu temporal
decay}\label{omoriutsu-temporal-decay}

Where a system has a meaningful event time series, we estimate temporal
aftershock decay following the Omori--Utsu form n(t) = K / (t + c)ᵖ
{[}6{]}. We identify a main-shock threshold by percentile or by a
multiple of the standard deviation, stack post-trigger event counts
across all main shocks in a forward window, log-bin the stack, and fit
(p, c, K) by weighted log-log regression with c grid-searched. Goodness
of fit is reported as a weighted R² in log space. Phase 4's
avalanche-size observables are fitted for size scaling only; the Omori
stack is not applied where the class makes no temporal-relaxation
prediction.

\subsubsection{2.5 Synthetic null
controls}\label{synthetic-null-controls}

For each phase we generate matched-\texttt{n} synthetic samples from
non-power-law sources and run the identical pipeline on each. Passing
requires correct rejection: the synthetic-null likelihood-ratio against
the matching alternative must be strongly negative, or the fit must fail
to converge on a stable \texttt{x\_min}. Phase 5 is a dedicated
null-control phase across four canonical non-SOC sources (Section 3.5).

\begin{center}\rule{0.5\linewidth}{0.5pt}\end{center}

\section{Five validation phases}\label{sec:3}

\subsubsection{3.1 Phase 1 --- USGS earthquakes (ground-truth gate
test)}\label{phase-1-usgs-earthquakes-ground-truth-gate-test}

\textbf{Data.} 84,724 tectonic earthquakes from the USGS Federated
Digital Seismograph Network (FDSN) event service, 2020-01-01 to
2025-01-01, M ≥ 3.5, restricted to \texttt{type\ =\ earthquake}. No
declustering is applied at catalog construction time, since declustering
would bias the Gutenberg--Richter fit being measured.

\textbf{Method and result.} The magnitude of completeness, estimated by
the Wiemer--Wyss maximum-curvature method, is M\_c = 4.45, leaving
37,281 events above completeness. An Aki maximum-likelihood fit with the
Shi--Bolt uncertainty yields a Gutenberg--Richter b-value of
\textbf{1.084 ± 0.005}, with a 500-resample bootstrap 95 \% confidence
interval of {[}1.073, 1.094{]}, implying an energy power-law exponent τ
= 1 + b/1.5 = 1.722. An independent Clauset--Shalizi--Newman
continuous-power-law fit on seismic energies s = 10¹·⁵ᴹ recovers α =
1.794 ± 0.024 (n = 1,071 above \texttt{x\_min}), consistent with the
b/1.5 relation under Hanks--Kanamori scaling. For aftershock sequences
following 580 main shocks of M ≥ 6.0 (24,680 stacked aftershocks at M ≥
4.0), a log-binned Omori--Utsu fit gives \textbf{p = 0.941 ± 0.017}, c =
0.10 d, weighted R² = 0.9927 over three temporal decades.

\textbf{Role.} Phase 1 is the gate test. Before running the pipeline on
any non-physics target, it must recover canonical behavior on a system
whose ground truth is not in dispute. It does: both exponents fall
inside canonical seismological ranges.

\subsubsection{3.2 Phase 2 --- S\&P 500 daily returns (first
cross-domain
transfer)}\label{phase-2-sp-500-daily-returns-first-cross-domain-transfer}

\textbf{Data.} 9,066 daily close prices of the S\&P 500 index
(\texttt{\^{}GSPC}, Yahoo Finance via the \texttt{yfinance} package,
1990-01-01 to 2025-12-31), giving 9,065 log returns with σ = 0.0114 and
range {[}−0.1277, +0.1096{]}.

\textbf{Method and result.} The same Clauset continuous power-law fit,
applied to the unsigned daily returns \textbar r\textbar, returns
\textbf{α = 2.998 ± 0.041} on 9,060 returns (n\_tail = 2,327 above
\texttt{x\_min} = 0.00998, a ≈1 \% daily move) --- reproducing the
Gopikrishnan et al.~``inverse cubic law'' to within 0.07 \% of the
canonical value of 3.0. An Omori--Utsu fit on stacked post-shock
volatility from 318 main shocks (\textbar r\textbar{} \textgreater{} 3σ,
threshold ≈ 3.42 \%) yields \textbf{p = 0.286 ± 0.034} (R² = 0.71),
inside the published daily-scale band of {[}0.3, 0.6{]} but outside the
intraday band of {[}0.7, 1.0{]} --- a scale-dependent feature, not a
deviation.

\textbf{Lognormal comparison --- direction stated correctly.} Table 1 of
the arxiv-02 source paper reports, for the power-law-vs-lognormal
likelihood-ratio test on these data, the normalized log-likelihood ratio
\textbf{R = −6.12} with significance \textbf{p = 9.3 × 10⁻¹⁰}. We adopt
the Clauset--Shalizi--Newman 2009 convention (their Eq. C.5 and §C; see
also Vuong 1989) in which the test statistic R is the sum of pointwise
log-likelihood differences and is \emph{positive when the power-law is
favored and negative when the alternative is favored}; under that
convention an R many standard deviations below zero with a small
two-sided p-value is a statistically distinguishable preference
\emph{for the alternative}. The arxiv-02 abstract and §3.1 prose (``the
power-law strongly dominating lognormal, p \textless{} 10⁻⁹'') is
therefore not a numerical mistake --- the underlying R = −6.12 is
correctly reported in its Table 1 --- but a \emph{sign-interpretation
error}: the small p-value is read as evidence for the power-law without
inspecting which direction R has fallen in. The corrected reading is
that for unsigned S\&P 500 daily returns the raw-tail Vuong test
\textbf{favors lognormal over the power-law at the \textbar R\textbar{}
= 6.12 level} (≈ 6σ on the standardized test statistic). The
thirteen-system sibling manuscript reports R = −6.12 with the same sign
convention and reads it correctly. We therefore restate the corrected
direction here. The SOC verdict for S\&P 500 (Section 4) does not rest
on rejecting the lognormal; it rests on \textbf{exponent-band agreement}
--- the measured α = 2.998 ± 0.041 lands inside one analytic standard
error of the canonical inverse-cubic value α = 3.0 {[}Gopikrishnan et
al.~1998, ref 21{]} --- and on the joint signature (power-law functional
form, Omori decay at the daily band, null controls passing). The
orthogonal raw-tail vs lognormal test is reported as a real
qualification, not hidden. The power-law-vs-exponential test for the
same fit is inconclusive (R = −0.52, p = 0.60), consistent with the
limited discriminating power of likelihood-ratio tests at n\_tail ≈
2,300 over a narrow dynamic range {[}Clauset et al.~2009 §6.3, ref
10{]}. The general implication of this --- that a raw-tail lognormal
comparison can favor lognormal without falsifying SOC membership --- is
discussed in §6.1.

\emph{Important scope qualification (added at v0.3 reviewer pre-empt).}
The canonical ``inverse cubic law'' α ≈ 3 is the \emph{empirically
observed} tail exponent of stock-return distributions {[}Gopikrishnan et
al.~1998, Plerou et al.~1999, Gabaix et al.~2003{]}, not an SOC
universality class exponent derived from first principles. The strongest
published form of the law is on \emph{individual-stock} and
\emph{sub-daily-resolution} data, where n\_tail reaches 10⁵--10⁶ and the
Vuong test has good discriminating power. The C1 measurement is the
daily-index transfer of the law (S\&P 500 daily, n\_tail = 2,327), which
has independent published confirmation in Weber et al.~2007 (ref 26) and
Petersen et al.~2010 (ref 25) but is the weaker form of the claim. The
result here is ``the S\&P 500 daily return tail reproduces the empirical
inverse-cubic value to within one analytic standard error,'' not ``the
S\&P 500 reproduces the SOC-predicted exponent.''

\subsubsection{3.3 Phase 3 --- DeFi liquidation cascades across three
protocols}\label{phase-3-defi-liquidation-cascades-across-three-protocols}

\textbf{Data.} 43,065 on-chain liquidation events from three
architecturally distinct DeFi lending protocols: Aave V2 (auction-based;
25,601 stablecoin-debt events), Compound V2 (direct liquidation with
incentive spread; 11,244 stablecoin-debt events), and MakerDAO's
Dog/Clip Liquidation 2.0 (Dutch clipper auctions; 1,985 events). Block
ranges span 2020 to 2024.

\textbf{Method and result.} The same Clauset fit gives tail exponents
\textbf{α = 1.684 ± 0.010} (Aave), \textbf{1.649 ± 0.016} (Compound),
and \textbf{1.567 ± 0.015} (MakerDAO) --- a spread of only 0.12 across
three independent codebases and liquidation mechanisms. Omori--Utsu
decay at 1-hour aggregation gives \textbf{p = 0.733 ± 0.045, 0.761 ±
0.042, 0.692 ± 0.071} respectively (weighted R² ∈ {[}0.24, 0.36{]},
lower than Phase 1 because DeFi event rates are sparse per hourly bin;
the slope is nonetheless far from zero --- at Compound the flat-rate
null is rejected at ≈18σ). Every per-protocol power-law fit decisively
rejects both lognormal and exponential alternatives (p \textless{}
10⁻⁹).

\textbf{Interpretation.} Three different liquidation mechanisms
producing α values within 0.12 of each other is the cross-instance
consistency a universality-class claim requires. With per-protocol σ(α)
≈ 0.010--0.016 the 0.12 spread is technically 7--12 standard errors ---
the three α values are \emph{not} statistically identical, and the
source paper does not claim they are --- but the spread is small on any
practical scale and well separated from the stock-return regime (α ≈
3.0). The DeFi exponents form a tight sub-cluster near earthquake energy
exponents (α ≈ 1.6--1.8), evidence for a ``discrete threshold-cascade''
sub-class of SOC spanning geology and decentralized finance.

\subsubsection{3.4 Phase 4 --- neural avalanches on task-active mouse
cortex}\label{phase-4-neural-avalanches-on-task-active-mouse-cortex}

\textbf{Data.} DANDI Archive dataset 000006, session
\texttt{sub-anm369962\_ses-20170313} --- 1,392,414 spikes from 71 sorted
units recorded over a 2,266 s delay-response behavioral task in mouse
anterior lateral motor (ALM) cortex. A synthetic check uses 200,000
critical Bienaymé--Galton--Watson branching-process avalanches.

\textbf{Method and result.} On the synthetic generator the pipeline
recovers the mean-field predictions cleanly: τ = 1.497 (predicted
1.500), α\_T = 1.917 (predicted 2.000), with the power-law form
dominating lognormal (p = 7 × 10⁻¹⁶) and exponential for both P(s) and
P(T). On the real recording, a bin-factor sweep across 1× to 16× the
mean inter-event interval gives two findings. First, the SOC scaling
relation γ = (α\_T − 1)/(τ − 1) is satisfied to within 2 \% at every bin
scale (measured \textbf{γ ≈ 1.10}, with the weighted-regression fit at
R² ≈ 0.998--0.999), and its stability across a 16-fold binning range is
a strong statistical signature of genuine criticality. Second, the
specific exponents are \textbf{not} the canonical mean-field
Beggs--Plenz values: \textbf{τ ∈ {[}2.17, 3.00{]}} and α\_T ∈ {[}2.49,
2.94{]} depending on bin width, consistently larger than τ = 3/2, α\_T =
2.

\textbf{Interpretation.} This is not a failure of criticality.
Task-active and subsampled cortical recordings are known to shift
exponents upward (Priesemann et al.); the recording therefore sits in a
task-active SOC sub-class rather than the spontaneous-activity
mean-field regime. The equivalence-class claim that neural avalanches
belong with earthquakes and DeFi liquidations is not contradicted --- it
is refined by the observation that the sub-class depends on brain state.

\subsubsection{3.5 Phase 5 --- synthetic non-SOC null
controls}\label{phase-5-synthetic-non-soc-null-controls}

\textbf{Purpose.} Phases 1--4 each found power-law tails. A standard
reviewer concern is ``would the pipeline fit a power law to noise too?''
A positive null result would fatally weaken every prior claim.

\textbf{Data and result.} The identical pipeline is run on four
synthetic datasets with known non-SOC distributions. As noted in §2.3, a
\emph{negative} R favors the alternative --- so for these nulls, a
strongly negative R is the \emph{correct} (passing) outcome:

{\def\LTcaptype{none} % do not increment counter
\begin{longtable}[]{@{}
  >{\raggedright\arraybackslash}p{(\linewidth - 8\tabcolsep) * \real{0.2000}}
  >{\raggedright\arraybackslash}p{(\linewidth - 8\tabcolsep) * \real{0.2000}}
  >{\raggedright\arraybackslash}p{(\linewidth - 8\tabcolsep) * \real{0.2000}}
  >{\raggedright\arraybackslash}p{(\linewidth - 8\tabcolsep) * \real{0.2000}}
  >{\raggedright\arraybackslash}p{(\linewidth - 8\tabcolsep) * \real{0.2000}}@{}}
\toprule\noalign{}
\begin{minipage}[b]{\linewidth}\raggedright
Null source
\end{minipage} & \begin{minipage}[b]{\linewidth}\raggedright
n
\end{minipage} & \begin{minipage}[b]{\linewidth}\raggedright
Fitted ``α''
\end{minipage} & \begin{minipage}[b]{\linewidth}\raggedright
Likelihood ratio
\end{minipage} & \begin{minipage}[b]{\linewidth}\raggedright
Verdict
\end{minipage} \\
\midrule\noalign{}
\endhead
\bottomrule\noalign{}
\endlastfoot
Gaussian random-walk increments (folded normal) & 20,000 & 2.999 & R =
−28.58 (vs lognormal), −44.76 (vs exponential) & power-law
\textbf{rejected} ✅ \\
Exponential variates & 20,000 & 2.996 & R = −16.03 (vs lognormal),
−17.17 (vs exponential) & power-law \textbf{rejected} ✅ \\
Homogeneous Poisson inter-arrival times & \textasciitilde50,000 & 3.000
& R = −24.45 (vs lognormal), −24.39 (vs exponential) & power-law
\textbf{rejected} ✅ \\
Homogeneous Poisson → Omori stack & 5,006 s window & --- & Omori p =
−0.068 (wrong sign), R² = 0.0015 & no Omori structure ✅ \\
\end{longtable}
}

Across three independent non-power-law size distributions the pipeline
correctly rejected the power-law hypothesis at likelihood ratios of −16
to −45; on the temporal side, the Omori detector on a homogeneous
Poisson process gave R² ≈ 0.002. By contrast, the real-data Phases 1--4
returned likelihood ratios favoring power-law and Omori R² values of
0.24 to 0.99.

\textbf{Conclusion.} The pipeline is a meaningful detector. It does not
confound genuine heavy-tailed SOC behavior with noise or with
exponential tails, and the positive findings of Phases 1--4 cannot be
dismissed as a methodological artifact. Null validations of power-law
fitting tools are an established practice {[}9{]}; Phase 5 is the
application of that standard practice to this specific pipeline.

\begin{center}\rule{0.5\linewidth}{0.5pt}\end{center}

\section{Cross-domain comparison}\label{sec:4}

Table 1 places the four real systems side by side. The headline
observation is that one fixed pipeline, applied with zero per-domain
re-tuning, recovers a coherent SOC signature in each of geophysics,
equity finance, decentralized finance, and neuroscience, while Phase 5
confirms the pipeline does not manufacture that signature from non-SOC
data.

\textbf{Table 1.} Four-system summary. Omori p is reported where the
system has a meaningful event time series.

{\def\LTcaptype{none} % do not increment counter
\begin{longtable}[]{@{}
  >{\raggedright\arraybackslash}p{(\linewidth - 10\tabcolsep) * \real{0.1667}}
  >{\raggedright\arraybackslash}p{(\linewidth - 10\tabcolsep) * \real{0.1667}}
  >{\raggedright\arraybackslash}p{(\linewidth - 10\tabcolsep) * \real{0.1667}}
  >{\raggedright\arraybackslash}p{(\linewidth - 10\tabcolsep) * \real{0.1667}}
  >{\raggedright\arraybackslash}p{(\linewidth - 10\tabcolsep) * \real{0.1667}}
  >{\raggedright\arraybackslash}p{(\linewidth - 10\tabcolsep) * \real{0.1667}}@{}}
\toprule\noalign{}
\begin{minipage}[b]{\linewidth}\raggedright
Phase
\end{minipage} & \begin{minipage}[b]{\linewidth}\raggedright
System
\end{minipage} & \begin{minipage}[b]{\linewidth}\raggedright
n (analysis sample)
\end{minipage} & \begin{minipage}[b]{\linewidth}\raggedright
Tail exponent
\end{minipage} & \begin{minipage}[b]{\linewidth}\raggedright
Omori p
\end{minipage} & \begin{minipage}[b]{\linewidth}\raggedright
Verdict
\end{minipage} \\
\midrule\noalign{}
\endhead
\bottomrule\noalign{}
\endlastfoot
1 & USGS earthquakes & 37,281 above M\_c & b = 1.084 ± 0.005 (α\_E =
1.794 ± 0.024) & 0.941 ± 0.017 & confirmed (ground-truth gate) \\
2 & S\&P 500 daily returns & 9,060 & α = 2.998 ± 0.041 & 0.286 ± 0.034 &
confirmed on exponent band (raw-tail LR favors lognormal --- see §3.2,
§6.1) \\
3a & Aave V2 liquidations & 25,601 & α = 1.684 ± 0.010 & 0.733 ± 0.045 &
confirmed \\
3b & Compound V2 liquidations & 11,244 & α = 1.649 ± 0.016 & 0.761 ±
0.042 & confirmed \\
3c & MakerDAO Dog liquidations & 1,985 & α = 1.567 ± 0.015 & 0.692 ±
0.071 & confirmed \\
4 & Mouse ALM cortex avalanches & 1,392,414 spikes & τ ∈ {[}2.17,
3.00{]} & --- (size-scaling phase) & sub-class shift; scaling relation γ
≈ 1.10 holds \\
5 & Synthetic non-SOC nulls (×4) & 20,000 each & rejected & rejected (R²
≈ 0.002) & correctly negative \\
\end{longtable}
}

The exponent spread across real systems --- α ≈ 1.6 (DeFi, earthquake
energy) to α ≈ 3.0 (S\&P 500) --- is \emph{expected} under
universality-class theory. Systems in the same class share equations of
motion, not necessarily a single numerical exponent: different conjugate
observables (released energy, return magnitude, debt size, avalanche
size) carry different scaling exponents. What the framework predicts to
be shared is the \emph{functional form} --- a power-law tail with an
exponential finite-size cutoff --- and the \emph{paired-signature
structure} (size power-law plus Omori temporal decay) for the
threshold-cascade members.

The most striking single number is the within-Phase-3 consistency: three
DeFi protocols with completely different liquidation engines (English
auction, direct incentive spread, Dutch clipper) yielding tail exponents
within 0.12 of each other. That is the kind of mechanism-independent
quantitative agreement that distinguishes a structural universality
claim from a surface analogy.

\begin{center}\rule{0.5\linewidth}{0.5pt}\end{center}

\section{Discussion}\label{sec:5}

\textbf{One pipeline, four domains.} The central claim of this preprint
is methodological as much as physical. Each of the four real systems has
been studied before, individually, in its own literature ---
Gutenberg--Richter for earthquakes, the inverse cubic law for equity
returns, branching-process criticality for neural avalanches. What has
not been done at this breadth is to fix a single Clauset-grade fitting
stack and transfer it, with no re-tuning, across geophysics, equity
finance, decentralized finance, and neuroscience. The fact that the same
code path recovers a canonical signature in each domain is the empirical
content of the ``they belong to one universality class'' claim that the
project's Layer 2 community-discovery step proposed from mechanism
graphs alone.

\textbf{Why the null phase is load-bearing.} A cross-domain power-law
survey is only persuasive if the surveying instrument can also say
``no.'' Phase 5 is therefore not an appendix but a core result: it
converts ``the pipeline found power laws everywhere it looked'' from a
worrying observation into a meaningful one, because the pipeline
demonstrably does \emph{not} find power laws in folded-normal,
exponential, or Poisson data, and does not find Omori decay in a
homogeneous Poisson process.

\textbf{DeFi as the flagship transfer.} Of the four real systems, DeFi
liquidations are the one with no prior published scaling measurement;
earthquakes, equity returns, and neural avalanches all have established
literature exponents to recover. Phase 3 is thus the phase where the
pipeline makes a genuinely new measurement rather than reproducing a
known one, and the three-protocol consistency is what gives that new
measurement its credibility.

\textbf{Refinement, not contradiction, in Phase 4.} The neural phase
deserves emphasis because it shows the framework behaving like a real
scientific instrument rather than a confirmation machine. The
recording's exponents are \emph{not} the textbook mean-field values; a
naïve ``confirm SOC'' pipeline would either force-fit them or report a
failure. Instead, the scaling relation γ ≈ 1.10 holds robustly across a
16-fold binning range --- the criticality signature that is invariant to
bin choice --- while the specific exponents correctly identify a
task-active sub-class. The class membership is refined, not rejected.

\begin{center}\rule{0.5\linewidth}{0.5pt}\end{center}

\section{Limitations}\label{sec:6}

This section is deliberately explicit. The honest scope of the result is
narrower than ``we proved cross-domain SOC.''

\textbf{6.1 Lognormal is not always rejected --- and is favored for S\&P
500.} The hardest alternative to a power-law tail is a lognormal, which
can mimic power-law behavior over a finite dynamic range. On the
raw-tail likelihood-ratio test, the lognormal alternative is not
decisively rejected for every system; for S\&P 500 daily returns it is
in fact \emph{favored} (R = −6.12, p = 9.3 × 10⁻¹⁰; see §3.2). The SOC
verdict for that system rests on exponent-band agreement --- α = 2.998
essentially on the canonical inverse-cubic 3.0 --- not on rejecting the
lognormal. We note this is consistent with a known procedural tension:
the raw-tail Vuong likelihood ratio (sensitive to a small number of
upper-tail individual events, where a lognormal can have a deeper
finite-n tail) and log-binned model selection can disagree, and at tail
sizes n\_tail ∼ 10³--10⁴ the power-law-vs-lognormal distinction is
partly procedural. The thirteen-system sibling manuscript discusses this
at length and reports that, on a complementary log-binned
Bayesian-information-criterion test, power-law-with-cutoff is preferred
over lognormal in every system tested. \textbf{The defensible framing is
that the SOC verdict rests on the joint signature --- power-law
functional form, exponents inside predicted bands, paired Omori decay,
null controls passing --- not on any single likelihood-ratio test.} The
raw-tail lognormal result for S\&P 500 is reported here as a real
qualification, not hidden.

\textbf{6.2 The pipeline detects endogenous threshold-cascade signatures
only.} This pipeline tests for the SOC threshold-cascade signature:
self-organized, slowly driven, internally generated cascades. It is not
constructed to detect, and should not be expected to flag, externally
driven or exogenous crises. We state this as a \emph{known structural
property of SOC threshold-cascade analysis methods in general} --- the
SOC model class describes endogenously organized criticality, and a
detector tuned to its power-law-plus-Omori signature has no sensitivity
to a shock imposed from outside the system {[}3, 4; and the SOC
monograph literature, e.g.~Jensen 1998, Pruessner 2012, Turcotte
1999{]}. It is \textbf{not} a finding of the present preprint, and the
Phase 1--5 source papers contain no within-project experiment that
demonstrates an endogenous/exogenous discrimination. Phase 5 validates
the pipeline only in the null-rejection direction (it correctly says
``no'' to non-SOC synthetic data); it does not establish that the
pipeline would catch every real-world crisis, and externally driven
events are explicitly outside the class the pipeline targets. A reviewer
extending the discussion of this point should cite the external SOC
literature, not this preprint.

\textbf{6.3 The project's downstream predictions are unverified.} The
Structural Isomorphism project's Layer 4 issues falsifiable numerical
predictions for systems that have \emph{not yet been measured}. The
current calibration status is honest about this: there are 24
predictions across 21 candidate classes, all with status ``to-validate''
(pending verification). For these predictions the target data has not
been collected --- the data sources are external databases --- so
\textbf{there is no observed sample to bootstrap a frequentist
confidence interval against}. The intervals attached to those
predictions are prior-based credible intervals, not confidence intervals
on data. \textbf{This preprint therefore makes no claim that the
project's cross-domain predictions are validated; only that the four
real systems above were measured with one pipeline and the null phase
passed.}

\textbf{6.4 Small-n and uncertainty caveats.} The MakerDAO sub-phase
(1,985 events) and the Phase 4 small-bin avalanche fits at 8×/16×
binning (n ≈ 19,000--38,000 avalanches, σ(τ) up to ≈0.11) are at the
lower end of where the Clauset likelihood-ratio test has good power;
``inconclusive'' results in those regimes should not be over-read. The
uncertainty quantification is not uniform across phases: Phase 1 reports
a 500-resample bootstrap on the b-value, whereas Phases 2--4 report the
analytic Clauset/Hill standard error (see §2.2 and Appendix Table A2). A
uniform bootstrap report across all four real phases would be a
reasonable methodological upgrade but was not done in the source papers.

\textbf{6.5 Single sessions / single windows.} Phase 4 rests on one
mouse-cortex recording session; the task-active-sub-class conclusion
would be strengthened by additional sessions. Phase 2 uses one index
over one time window. None of the phases is a meta-analysis across many
independent datasets within the same domain.

\textbf{6.6 This is a Claude-generated draft.} Every number in this
draft traces to one of the five Phase 1--5 source papers (cross-checked
against the thirteen-system sibling manuscript), but the synthesis,
framing, and cross-phase claims have not been checked by a domain
expert. The Pre-submission checklist at the end of this file lists every
remaining human-confirmation item.

\begin{center}\rule{0.5\linewidth}{0.5pt}\end{center}

\section{Conclusion}\label{sec:7}

We assembled a single Clauset-grade analysis pipeline and applied it,
without per-domain tuning, to four independent real systems --- USGS
earthquakes, S\&P 500 daily returns, DeFi liquidations across three
protocols, and task-active mouse-cortex neural avalanches --- plus four
synthetic non-SOC null sources. The pipeline recovered canonical SOC
signatures on all four real systems (Gutenberg--Richter b = 1.084 ±
0.005; inverse-cubic α = 2.998 ± 0.041; DeFi α ∈ {[}1.567, 1.684{]} with
cross-protocol spread 0.12; neural scaling relation γ ≈ 1.10 stable
across a 16-fold binning range) and correctly rejected the power-law
hypothesis on all four nulls. The result is a single-pipeline,
cross-domain test of self-organized-criticality universality,
deliberately conservative in its claims: the lognormal alternative is
not rejected in every raw-tail test --- for S\&P 500 it is favored, and
the SOC verdict there rests on exponent-band agreement --- the pipeline
targets only endogenous threshold-cascade dynamics, and the project's
downstream cross-domain predictions remain unverified for lack of target
data. Within those limits, four very different systems gave four
coherent results from one method --- which is the minimum empirical bar
a universality-class claim must clear.

\begin{center}\rule{0.5\linewidth}{0.5pt}\end{center}

\section*{References}\label{references}

The list below is consolidated and de-duplicated from the individually
checked bibliographies of the five Phase 1--5 source papers. Each entry
was cross-checked against at least one source paper's reference list.
Entries that could not be cross-verified within the project repository
are marked \textbf{[to verify]} (verify bibliographic detail before
submission).

\textbf{Foundational --- universality and SOC}

\begin{enumerate}
\def\labelenumi{\arabic{enumi}.}
\tightlist
\item
  Wilson, K. G. The renormalization group and critical phenomena.
  \emph{Reviews of Modern Physics} \textbf{55}, 583 (1983).
\item
  Stanley, H. E. Scaling, universality, and renormalization: three
  pillars of modern critical phenomena. \emph{Reviews of Modern Physics}
  \textbf{71}, S358 (1999).
\item
  Bak, P., Tang, C. \& Wiesenfeld, K. Self-organized criticality: an
  explanation of 1/f noise. \emph{Physical Review Letters} \textbf{59},
  381 (1987).
\item
  Olami, Z., Feder, H. J. S. \& Christensen, K. Self-organized
  criticality in a continuous, nonconservative cellular automaton
  modeling earthquakes. \emph{Physical Review Letters} \textbf{68}, 1244
  (1992).
\item
  Turcotte, D. L. Self-organized criticality. \emph{Reports on Progress
  in Physics} \textbf{62}, 1377 (1999).
\item
  Sethna, J. P., Dahmen, K. A. \& Myers, C. R. Crackling noise.
  \emph{Nature} \textbf{410}, 242 (2001).
\item
  Jensen, H. J. \emph{Self-Organized Criticality: Emergent Complex
  Behavior in Physical and Biological Systems.} Cambridge University
  Press (1998).
\item
  Pruessner, G. \emph{Self-Organised Criticality: Theory, Models and
  Characterisation.} Cambridge University Press (2012).
\item
  Sornette, D. \emph{Critical Phenomena in Natural Sciences: Chaos,
  Fractals, Selforganization and Disorder.} 2nd ed., Springer (2006).
\end{enumerate}

\textbf{Statistical method --- power-law fitting}

\begin{enumerate}
\def\labelenumi{\arabic{enumi}.}
\setcounter{enumi}{9}
\tightlist
\item
  Clauset, A., Shalizi, C. R. \& Newman, M. E. J. Power-law
  distributions in empirical data. \emph{SIAM Review} \textbf{51}, 661
  (2009).
\item
  Alstott, J., Bullmore, E. \& Plenz, D. powerlaw: a Python package for
  analysis of heavy-tailed distributions. \emph{PLoS ONE} \textbf{9},
  e85777 (2014).
\item
  Newman, M. E. J. Power laws, Pareto distributions and Zipf's law.
  \emph{Contemporary Physics} \textbf{46}, 323 (2005).
\item
  Broido, A. D. \& Clauset, A. Scale-free networks are rare.
  \emph{Nature Communications} \textbf{10}, 1017 (2019).
\end{enumerate}

\textbf{Earthquakes (Phase 1)}

\begin{enumerate}
\def\labelenumi{\arabic{enumi}.}
\setcounter{enumi}{13}
\tightlist
\item
  Gutenberg, B. \& Richter, C. F. Frequency of earthquakes in
  California. \emph{Bulletin of the Seismological Society of America}
  \textbf{34}, 185 (1944).
\item
  Omori, F. On the after-shocks of earthquakes. \emph{Journal of the
  College of Science, Imperial University of Tokyo} \textbf{7}, 111
  (1894).
\item
  Utsu, T., Ogata, Y. \& Matsu'ura, R. S. The centenary of the Omori
  formula for a decay law of aftershock activity. \emph{Journal of
  Physics of the Earth} \textbf{43}, 1 (1995).
\item
  Aki, K. Maximum likelihood estimate of b in the formula log N = a − bM
  and its confidence limits. \emph{Bulletin of the Earthquake Research
  Institute, University of Tokyo} \textbf{43}, 237 (1965).
\item
  Shi, Y. \& Bolt, B. A. The standard error of the magnitude-frequency b
  value. \emph{Bulletin of the Seismological Society of America}
  \textbf{72}, 1677 (1982).
\item
  Wiemer, S. \& Wyss, M. Minimum magnitude of completeness in earthquake
  catalogs: examples from Alaska, the western United States, and Japan.
  \emph{Bulletin of the Seismological Society of America} \textbf{90},
  859 (2000).
\item
  Hanks, T. C. \& Kanamori, H. A moment magnitude scale. \emph{Journal
  of Geophysical Research} \textbf{84}, 2348 (1979).
\end{enumerate}

\textbf{Equity finance (Phase 2)}

\begin{enumerate}
\def\labelenumi{\arabic{enumi}.}
\setcounter{enumi}{20}
\tightlist
\item
  Gopikrishnan, P., Meyer, M., Amaral, L. A. N. \& Stanley, H. E.
  Inverse cubic law for the distribution of stock price variations.
  \emph{European Physical Journal B} \textbf{3}, 139 (1998).
\item
  Plerou, V., Gopikrishnan, P., Amaral, L. A. N., Meyer, M. \& Stanley,
  H. E. Scaling of the distribution of price variations of individual
  companies. \emph{Physical Review E} \textbf{60}, 6519 (1999).
\item
  Gabaix, X., Gopikrishnan, P., Plerou, V. \& Stanley, H. E. A theory of
  power-law distributions in financial market fluctuations.
  \emph{Nature} \textbf{423}, 267 (2003).
\item
  Lillo, F. \& Mantegna, R. N. Power-law relaxation in a complex system:
  Omori law after a financial market crash. \emph{Physical Review E}
  \textbf{68}, 016119 (2003).
\item
  Petersen, A. M., Wang, F., Havlin, S. \& Stanley, H. E. Market
  dynamics immediately before and after financial shocks: quantifying
  the Omori, productivity, and Bath laws. \emph{Physical Review E}
  \textbf{82}, 036114 (2010).
\item
  Weber, P., Wang, F., Vodenska-Chitkushev, I., Havlin, S. \& Stanley,
  H. E. Relation between volatility correlations in financial markets
  and Omori processes occurring on all scales. \emph{Physical Review E}
  \textbf{76}, 016109 (2007).
\item
  Mantegna, R. N. \& Stanley, H. E. \emph{An Introduction to
  Econophysics: Correlations and Complexity in Finance.} Cambridge
  University Press (2000).
\end{enumerate}

\textbf{DeFi (Phase 3)}

\begin{enumerate}
\def\labelenumi{\arabic{enumi}.}
\setcounter{enumi}{27}
\tightlist
\item
  Qin, K., Zhou, L. \& Gervais, A. An empirical study of DeFi
  liquidations: incentives, risks, and instabilities. \emph{Proceedings
  of the ACM Internet Measurement Conference (IMC '21)} (2021).
\item
  Perez, D., Werner, S. M., Xu, J. \& Livshits, B. Liquidations: DeFi on
  a knife-edge. \emph{Financial Cryptography and Data Security 2021}
  (2021).
\item
  Aave. \emph{Aave Protocol V2 Whitepaper}. Technical documentation,
  Aave Companies (2020). URL:
  https://github.com/aave/aave-protocol/blob/master/docs/Aave\_Protocol\_Whitepaper\_v1\_0.pdf.
  Accessed: 2026-05-24. \emph{{[}Cited as technical documentation; no
  journal venue.{]}}
\item
  Compound Labs. \emph{Compound: The Money Market Protocol --- V2
  Whitepaper}. Technical documentation, Compound Labs (2019). URL:
  https://compound.finance/documents/Compound.Whitepaper.pdf. Accessed:
  2026-05-24. \emph{{[}Cited as technical documentation; no journal
  venue.{]}}
\item
  MakerDAO. \emph{Liquidation 2.0 (LIQ-2.0): Dog and Clipper
  specification.} Technical documentation, MakerDAO (2021). URL:
  https://docs.makerdao.com/smart-contract-modules/dog-and-clipper-detailed-documentation.
  Accessed: 2026-05-24. \emph{{[}Cited as online specification; no
  journal venue.{]}}
\item
  Motter, A. E. \& Lai, Y.-C. Cascade-based attacks on complex networks.
  \emph{Physical Review E} \textbf{66}, 065102 (2002).
\end{enumerate}

\textbf{Neural avalanches (Phase 4)}

\begin{enumerate}
\def\labelenumi{\arabic{enumi}.}
\setcounter{enumi}{33}
\tightlist
\item
  Beggs, J. M. \& Plenz, D. Neuronal avalanches in neocortical circuits.
  \emph{Journal of Neuroscience} \textbf{23}, 11167 (2003).
\item
  Friedman, N., Ito, S., Brinkman, B. A. W., Shimono, M., DeVille, R. E.
  L., Dahmen, K. A., Beggs, J. M. \& Butler, T. C. Universal critical
  dynamics in high resolution neuronal avalanche data. \emph{Physical
  Review Letters} \textbf{108}, 208102 (2012).
\item
  Priesemann, V., Munk, M. H. J. \& Wibral, M. Subsampling effects in
  neuronal avalanche distributions recorded in vivo. \emph{BMC
  Neuroscience} \textbf{15}, 1 (2014).
\item
  Touboul, J. \& Destexhe, A. Power-law statistics and universal scaling
  in the absence of criticality. \emph{Physical Review E} \textbf{95},
  012413 (2017).
\item
  Harris, T. E. \emph{The Theory of Branching Processes.} Springer
  (1963); Dover reprint (1989).
\item
  Li, N., Chen, T.-W., Guo, Z. V., Gerfen, C. R. \& Svoboda, K. A motor
  cortex circuit for motor planning and movement. \emph{Nature}
  \textbf{519}, 51 (2015). {[}Data: DANDI Archive dataset 000006, mouse
  anterior lateral motor cortex in delay-response task,
  https://dandiarchive.org/dandiset/000006.{]}
\item
  Beggs, J. M. \& Timme, N. Being critical of criticality in the brain.
  \emph{Frontiers in Physiology} \textbf{3}, 163 (2012).
\end{enumerate}

\textbf{Project (self-references)}

\emph{Reviewer note.} Refs 41--44 are companion arXiv preprints by the
same author covering Phases 1--4 individually. Each has been drafted
from the same Phase 1--5 source-paper materials used in this synthesis;
the arXiv identifiers will be filled in at the time C1 is submitted (the
four Phase papers and C1 are intended to be cross-linked by arXiv ID on
the same submission day). Ref 45 is the Structural Isomorphism Project's
Zenodo deposit --- see §1 of the Pre-submission checklist for the
deposit-coverage caveat (the DOI resolves to a Zenodo record but that
record at present covers the project's V1/V2 contrastive-learning
benchmark, \emph{not} the Phase 1--5 SOC code/data; a new Phase-1--5
deposit may be required before submission).

\begin{enumerate}
\def\labelenumi{\arabic{enumi}.}
\setcounter{enumi}{40}
\tightlist
\item
  Wan, Q. Recovering self-organized criticality on a global earthquake
  catalog: a reproducible pipeline for cross-domain universality-class
  identification. \emph{arXiv:2605.XXXXX} {[}physics.geo-ph{]} (2026).
  \emph{[arXiv ID to be filled at submission]}
\item
  Wan, Q. Cross-domain self-organized criticality: inverse cubic law and
  Omori decay on thirty-five years of S\&P 500 daily returns.
  \emph{arXiv:2605.XXXXX} {[}q-fin.ST{]} (2026). \emph{[arXiv ID to be filled at submission]}
\item
  Wan, Q. Cross-protocol SOC universality in DeFi liquidation cascades:
  43,065 events across Aave V2, Compound V2, and MakerDAO.
  \emph{arXiv:2605.XXXXX} {[}q-fin.ST{]} (2026). \emph{[arXiv ID to be filled at submission]}
\item
  Wan, Q. Criticality without mean-field SOC: neural avalanche scaling
  on task-active mouse cortex. \emph{arXiv:2605.XXXXX} {[}q-bio.NC{]}
  (2026). \emph{[arXiv ID to be filled at submission]}
\item
  Structural Isomorphism Project. Project snapshot: cross-domain
  universality-class identification (benchmark + code). Zenodo (2026),
  DOI: 10.5281/zenodo.19547879. \emph{{[}Reviewer note: this DOI
  currently resolves to the project's V1/V2 contrastive-learning
  benchmark, not to the Phase 1--5 SOC code/data --- see Pre-submission
  checklist item 1.{]}}
\end{enumerate}

\begin{center}\rule{0.5\linewidth}{0.5pt}\end{center}

\section*{Appendix A --- Data and reproducibility}\label{appendix-a}

\textbf{Table A1.} Data sources and analysis samples.

{\def\LTcaptype{none} % do not increment counter
\begin{longtable}[]{@{}
  >{\raggedright\arraybackslash}p{(\linewidth - 6\tabcolsep) * \real{0.2500}}
  >{\raggedright\arraybackslash}p{(\linewidth - 6\tabcolsep) * \real{0.2500}}
  >{\raggedright\arraybackslash}p{(\linewidth - 6\tabcolsep) * \real{0.2500}}
  >{\raggedright\arraybackslash}p{(\linewidth - 6\tabcolsep) * \real{0.2500}}@{}}
\toprule\noalign{}
\begin{minipage}[b]{\linewidth}\raggedright
Phase
\end{minipage} & \begin{minipage}[b]{\linewidth}\raggedright
System
\end{minipage} & \begin{minipage}[b]{\linewidth}\raggedright
Data source
\end{minipage} & \begin{minipage}[b]{\linewidth}\raggedright
Analysis sample
\end{minipage} \\
\midrule\noalign{}
\endhead
\bottomrule\noalign{}
\endlastfoot
1 & USGS earthquakes & USGS FDSN event service, 2020--2025, M ≥ 3.5 &
37,281 above M\_c = 4.45 \\
2 & S\&P 500 & Yahoo Finance \texttt{\^{}GSPC} daily close
(\texttt{yfinance}), 1990--2025 & 9,060 returns (n\_tail = 2,327) \\
3 & DeFi liquidations & On-chain event logs: Aave V2, Compound V2,
MakerDAO Dog/Clip & 43,065 events (25,601 / 11,244 / 1,985) \\
4 & Mouse cortex & DANDI Archive 000006, session
sub-anm369962\_ses-20170313 & 1,392,414 spikes / 71 units \\
5 & Synthetic nulls & Generated in-repo
(\texttt{v4/validation/null-controls}) & 20,000 each (×3 size, ×1
Omori) \\
\end{longtable}
}

\textbf{Table A2.} Uncertainty quantification per phase (closes the v0.1
§2.2 {[}TODO{]}).

{\def\LTcaptype{none} % do not increment counter
\begin{longtable}[]{@{}
  >{\raggedright\arraybackslash}p{(\linewidth - 6\tabcolsep) * \real{0.2500}}
  >{\raggedright\arraybackslash}p{(\linewidth - 6\tabcolsep) * \real{0.2500}}
  >{\raggedright\arraybackslash}p{(\linewidth - 6\tabcolsep) * \real{0.2500}}
  >{\raggedright\arraybackslash}p{(\linewidth - 6\tabcolsep) * \real{0.2500}}@{}}
\toprule\noalign{}
\begin{minipage}[b]{\linewidth}\raggedright
Phase
\end{minipage} & \begin{minipage}[b]{\linewidth}\raggedright
Quantity
\end{minipage} & \begin{minipage}[b]{\linewidth}\raggedright
Uncertainty method
\end{minipage} & \begin{minipage}[b]{\linewidth}\raggedright
Detail
\end{minipage} \\
\midrule\noalign{}
\endhead
\bottomrule\noalign{}
\endlastfoot
1 & b-value & 500-resample bootstrap \textbf{+} analytic Shi--Bolt error
& seed = 42; 95 \% CI {[}1.073, 1.094{]} \\
1 & α (seismic energy) & Clauset/Hill analytic standard error & ±
0.024 \\
2 & α (S\&P 500) & Clauset/Hill analytic standard error & ± 0.041; no
bootstrap in source paper \\
3 & α (each DeFi protocol) & Clauset/Hill analytic standard error & ±
0.010--0.016; no bootstrap in source paper \\
4 & τ, α\_T (per bin factor) & Clauset/Hill analytic standard error & ±
0.01--0.13 depending on bin; no bootstrap in source paper \\
\end{longtable}
}

\textbf{Pipeline.} Authoritative implementation: the
\texttt{soc-pipeline} Python package (\texttt{packages/soc-pipeline/},
version 0.1.0, MIT licence; depends on \texttt{numpy}, \texttt{scipy},
\texttt{pandas}, \texttt{powerlaw}). The legacy path
\texttt{v4/lib/soc\_pipeline.py} is a deprecation shim re-exporting the
package. Per-phase analysis scripts are tagged in the repository
(\texttt{v4/phase1-earthquake-2026-04-15},
\texttt{v4/phase2-stockmarket-2026-04-15}, \ldots). See §2 for the
provenance caveat regarding the ``339 lines / commit 7ee228c''
description in the thirteen-system sibling manuscript.

\textbf{Project Zenodo deposit.} DOI 10.5281/zenodo.19547879 (cited by
the Phase 1--4 source papers). \textbf{[pre-submission verification]} This DOI could not be
verified offline; see Pre-submission checklist item 1.

\begin{center}\rule{0.5\linewidth}{0.5pt}\end{center}

\section*{Pre-submission checklist}\label{checklist}

The seven v0.1 [TODO] markers have been closed in this v0.2
draft (see the META block changelog at the top). The items below are the
residual human-confirmation tasks that cannot be settled from the
repository alone:

\begin{enumerate}
\def\labelenumi{\arabic{enumi}.}
\tightlist
\item
  \textbf{Zenodo DOI.} Verify online that DOI
  \texttt{10.5281/zenodo.19547879} resolves, is the current version, and
  that its deposit covers all Phase 1--5 code and data referenced here.
  The DOI is shown in the text but flagged [to confirm]; do not present
  it as confirmed until checked.
\item
  \textbf{Pipeline canonical version.} Record one canonical release tag
  or commit hash of the \texttt{soc-pipeline} package against which the
  final Phase 1--5 numbers are confirmed, and reconcile it with the
  per-phase repository tags. Note for the reviewer: the thirteen-system
  sibling manuscript's ``\texttt{v4/lib/soc\_pipeline.py}, 339 lines,
  commit \texttt{7ee228c}'' description is stale relative to the current
  package-based layout.
\item
  \textbf{Reference entries marked [to verify].} Confirm full
  bibliographic detail for the entries flagged [to verify] in the
  reference list --- chiefly the DeFi whitepapers/specifications (refs
  30--32, which have no journal venue and should perhaps be cited as
  technical documentation with access dates) and the project
  self-references (refs 41--45, internal project documents whose final
  citable form / arXiv identifiers must be decided before submission).
\item
  \textbf{Phase 2 lognormal wording --- final editorial sign-off.}
  §3.2/§6.1 now state that the raw-tail Vuong test favors lognormal for
  S\&P 500 (R = −6.12) and that the arxiv-02 source paper's prose
  (``power-law strongly dominates lognormal'') is an internal
  sign-interpretation error. A domain reviewer should confirm this
  reading and decide whether to additionally issue a correction note for
  the standalone arxiv-02 paper.
\item
  \textbf{Decide whether to post the thirteen-system sibling alongside
  C1.} This is an editorial choice (post both, or only the focused
  five-system paper). The scope decision in §1 fixes C1 as the focused
  paper; whether the sibling is also posted is left to the author at
  submission time.
\item
  \textbf{Domain-expert review.} This entire draft is Claude-generated
  and has not been read by a domain expert (§6.6). A
  seismology/econophysics/neuroscience reviewer should check the
  cross-phase synthesis and framing before submission.
\end{enumerate}

\begin{center}\rule{0.5\linewidth}{0.5pt}\end{center}

\emph{End of draft v0.2. Status: draft, pending human review. Generated 2026-05-22 from
project source materials listed in the meta block at the top of this
file. The seven v0.1 {[}TODO{]} markers are closed; six residual
human-confirmation items are listed above. Not yet reviewed by a domain
expert.}

\end{document}
